\section{The \texttt{Brilliance\_monitor} McStas Component}
Special "Brilliance" monitor.

\subsection*{Identification}
\begin{itemize}
  \item \textbf{Author:} Peter Willendrup, derived from TOF\_lambda\_monitor.comp
  \item \textbf{Origin:} DTU Physics
  \item \textbf{Date:} May 23, 2012
\end{itemize}

\subsection*{Description}
If used in the right setting, will output "instantaneous" and "mean" brilliances in units of Neutrons/cm\textasciicircum{}2/ster/\AA{}/s. Conditions for proper units: \textless{}ul\textgreater{} \textless{}li\textgreater{}Use a with a source of area 1x1cm \textless{}li\textgreater{}The source must illuminate/focus to an area of 1x1cm a 1m distance \textless{}li\textgreater{}Parametrise the Brilliance\_monitor with the frequency of the source \textless{}li\textgreater{}To not change the source TOF distribution, place the Brilliance monitor close to the source! \textless{}/ul\textgreater{}

with a source of area 1x1cm illuminating/focusing to an area of 1x1cm a 1m distance, this monitor will output "instantaneous" and "mean" brilliances in units of Neutrons/cm\textasciicircum{}2/ster/\AA{}/s

Here is an example of the use of the component. Note how the mentioned Unit conditions are implemented in instrument code.

COMPONENT Source = ESS\_moderator\_long( l\_low = lambdamin, l\_high = lambdamax, dist = 1, xw = 0.01, yh = 0.01,

\begin{verbatim}
freq = 14, T=50, tau=287e-6, tau1=0, tau2=20e-6,
n=20, n2=5, d=0.00286, chi2=0.9, I0=6.9e11, I2=27.6e10,
\end{verbatim}

branch1=0, branch2=0.5, twopulses=0, size=0.01) AT (0, 0, 0) RELATIVE Origin

COMPONENT BRIL = Brilliance\_monitor(nlam=196,nt=401,filename="bril.sim", t\_0=0,t\_1=4000,lambda\_0=lambdamin, lambda\_1=lambdamax, Freq=14) AT (0,0,0.000001) RELATIVE Source

\subsection*{Input parameters}
Parameters in \textbf{boldface} are required; the others are optional.

\begin{longtable}{p{0.22\textwidth}p{0.12\textwidth}p{0.46\textwidth}p{0.14\textwidth}}
\toprule
\textbf{Name} & \textbf{Unit} & \textbf{Description} & \textbf{Default} \\
\midrule
\endhead
nlam & 1 & Number of bins in wavelength & 101 \\
nt & 1 & Number of bins in TOF & 1001 \\
nowritefile & 1 & If set, monitor will skip writing to disk & 0 \\
lambda\_0 & \AA{} & Minimum wavelength detected & 0 \\
lambda\_1 & \AA{} & Maximum wavelength detected & 20 \\
restore\_neutron & 1 & If set, the monitor does not influence the neutron state & 0 \\
\textbf{Freq} & Hz & Source frequency. Use freq=1 for reactor source &  \\
tofcuts & 1 & Flag to generate TOF-distributions as function of wavelength & 0 \\
toflambda & 1 & Flag to generate TOF-lambda distribution output & 0 \\
xwidth & m & width of monitor & 0.01 \\
yheight & m & height of monitor & 0.01 \\
source\_dist & m & Distance from source. Beware when transporting through neutron optics! & 1 \\
filename & string & Defines filenames for the detector images. Stored as:\textless{}br\textgreater{}Peak\_\&lt;filename\&gt; and Mean\_\&lt;filename\&gt; & 0 \\
t\_0 & us & Minimum time & 0 \\
t\_1 & us & Maximum time & 20000 \\
srcarea & cm$^{2}$ & Source area & 1 \\
\bottomrule
\end{longtable}

\subsection*{Links}
\begin{itemize}
  \item Component source code found in file \texttt{Brilliance\_monitor.comp}.
\end{itemize}
\IfFileExists{monitors/Brilliance_monitor_static.tex}{\input{monitors/Brilliance_monitor_static.tex}}{}